
\subsection{Author contributions}
O.N., E.B., L.E.W., J.H., C.C.H., A.J.A., I.E.B. and H.C. contributed to conceptualisation. O.N., E.B., M.F., L.E.W., C.C.H., A.J.A., I.E.B. and H.C. contributed to methodology. O.N., M.F. and H.C. developed the software. O.N., M.F., L.E.W. and H.C. performed formal analysis. O.N., E.B. and H.C. wrote the original draft. All authors contributed to manuscript review and editing. H.C. supervised the project.

\subsection{Funding}
The authors declare that this research did not receive any specific grant from funding agencies.

\subsection{Competing interests}
The authors declare no competing interests.

\subsection{Data availability}
Data access and analysis were conducted under the governance framework of the Cheshire \& Merseyside Data Into Action (DIA) programme. The study received approval from the Data Access \& Asset Group (DAAG) following completion of a Data Access Request Form (DARF) and Data Protection Impact Assessment (DPIA). Access was granted for the period June 8, 2025 to November 8, 2025, after which all access was revoked. The study used pseudonymised patient profiles and did not require individual patient consent under the approved data governance arrangements. Data controllers across NHS Cheshire \& Merseyside were informed of the study and retained the right to opt out their patient populations. We open source 45 detailed patient failure vignettes.